\section{引言}
本讲一开场便由 Xinyun Chen 提到“\textit{2025 will be the year of Agents}”，并说明本学期 Berkeley RDI 的 Advanced LLM Agents SP25 课程在 Google DeepMind 与 Cal 的协作下，既有应用 track 的成熟案例，也新增着眼于 inference-time reasoning 的研究 track。讲师强调 15,000+ 名 MOOC 报名者、3,000+ 名 hackathon 参与者以及超过 20 万美元的行业资源，意在表明本课程不是纯理论演示，而是真正要把推理时技术部署到复杂工程场景中。

\begin{knowledgebox}{课程定位}
模块化地把推理时技术拆成 “提示 + 搜索 + 迭代” 三部分，讲师将在本讲详细分享每一环的策略、指令模板与工程要点，让我们在“推理时投入算力”这个主题下同时把 token/width/depth 三个维度照顾到。
\end{knowledgebox}

\subsection{本章小结}
引言为本讲设定了 Agent 冷启动背景，强调 inference-time compute 作为不改权重即提升推理的杠杆，并把内容组织成提示、搜索与迭代三段，后续章节按照这一逻辑展开。

